% Generated by cv/build_cv.py from the selected profile and publications JSON.
% Layout matches the supplied academic CV; edit the JSON records, not this file.
\documentclass[letterpaper,10pt]{article}
\def\theme{MidnightBlue}
\renewcommand{\baselinestretch}{1.0}
\selectfont
\usepackage{fontspec}
\setmainfont{Times New Roman}
\PassOptionsToPackage{french}{babel}
\usepackage{simplecv}
\usepackage{needspace}

\begin{document}
\selectlanguage{french}
\headinginline{WEI Shiyi}{
    Courriel: \href{mailto:shiyi.wei@univ-cotedazur.fr}{shiyi.wei@univ-cotedazur.fr}
}

\section{Intérêts de recherche}

Physique statistique hors équilibre et systèmes désordonnés, diffusion quantique et transport mésoscopique, fluctuations de densité et de courant dans les systèmes quantiques bruités et grandes déviations et théorie des fluctuations macroscopiques.

\section{Formation}
\outerlist{
\begin{samepage}
\entrybigplus
{\textbf{Université Côte d’Azur}}{Nice, France}
{Doctorat en physique théorique}{Nov. 2024 -- oct. 2027 (prévu)}
{Directeur de thèse: Mathias Albert; Co-encadrant: Tony Jin}{}
\innerlist{
\entry{Projet de thèse initial : Towards a typical theory of quantum diffusion.}
}
\end{samepage}

\begin{samepage}
\entrybig
{\textbf{École Normale Supérieure}}{Paris, France}
{Diplôme de l’ENS, physique}{2020 -- 2024}
\innerlist{
\entry{Master 1 et Master 2 ICFP, parcours Physique quantique.}
}
\end{samepage}

\begin{samepage}
\entrybig
{\textbf{Sorbonne Université}}{Paris, France}
{Licence 1 et 2, double majeure Mathématiques et Physique}{2018 -- 2020}
\end{samepage}
}

\section{Publications}
\orderedouterlist{
\entry{Mathias Albert, Denis Bernard, Tony Jin, Stefano Scopa et \textbf{Shiyi Wei}. "Universal classical and quantum fluctuations in the large deviations of current of noisy quantum systems: Quantum symmetric simple exclusion and inclusion processes." \textit{Physical Review E} \textbf{114}, 034140 (2026). \href{https://doi.org/10.1103/dt4c-tz98}{doi:10.1103/dt4c-tz98}; \href{https://arxiv.org/abs/2601.16883}{arXiv:2601.16883}.}
\entry{Denis Bernard, Tony Jin, Stefano Scopa et \textbf{Shiyi Wei}. "Large deviations of density fluctuations in the boundary-driven quantum symmetric simple inclusion process." \textit{Physical Review E} \textbf{112}, 034106 (2025). \href{https://doi.org/10.1103/fs7f-z3k6}{doi:10.1103/fs7f-z3k6}; \href{https://arxiv.org/abs/2503.18763}{arXiv:2503.18763}.}
}

\Needspace{31\baselineskip}
\section{Expérience en recherche}
\outerlist{
\Needspace{8\baselineskip}
\begin{samepage}
\entrybig
{\parbox[t]{0.74\textwidth}{\raggedright\textbf{Institut de Physique de Nice (INPHYNI)}}}{Nice, France}
{Doctorante en physique théorique}{Nov. 2024 -- présent}
\par\vspace{2pt}{\small Encadrants: Mathias Albert et Tony Jin\par}
\innerlist{
\entry{Étude théorique et modélisation analytique et numérique des fluctuations dans les systèmes désordonnés hors équilibre.}
}
\end{samepage}

\Needspace{8\baselineskip}
\begin{samepage}
\entrybig
{\parbox[t]{0.74\textwidth}{\raggedright\textbf{Institut de Physique de Nice (INPHYNI)}}}{Nice, France}
{Stagiaire de Master 2}{Avr. -- juin 2024}
\par\vspace{2pt}{\small Encadrants: Mathias Albert et Tony Jin\par}
\innerlist{
\entry{Calcul analytique des cumulants de courant hors équilibre dans le processus quantique d’exclusion simple symétrique (QSSEP), à l’aide du formalisme de Lindblad et du calcul d’Itô quantique.}
}
\end{samepage}

\Needspace{8\baselineskip}
\begin{samepage}
\entrybig
{\parbox[t]{0.74\textwidth}{\raggedright\textbf{LPMMC (CNRS)}}}{Grenoble, France}
{Stagiaire de Master 1}{Fév. -- juil. 2022}
\par\vspace{2pt}{\small Encadrant: Anna Minguzzi\par}
\innerlist{
\entry{Modélisation d’un gaz de Tonks--Girardeau et simulation numérique de sa matrice densité avec Mathematica.}
}
\end{samepage}

\Needspace{8\baselineskip}
\begin{samepage}
\entrybig
{\parbox[t]{0.74\textwidth}{\raggedright\textbf{Laboratoire Kastler Brossel (LKB)}}}{Paris, France}
{Stagiaire de Licence 1}{Juin -- juil. 2019}
\par\vspace{2pt}{\small Encadrant: Alice Sinatra\par}
\innerlist{
\entry{Étude de la reconstruction d’états quantiques à l’aide des distributions de Wigner, des états cohérents et du squeezing quantique.}
}
\end{samepage}
}

\Needspace{8\baselineskip}
\section{Langues}
\outerlist{
\entry{\textbf{Français:} Courant}
\entry{\textbf{Anglais:} Courant}
\entry{\textbf{Chinois:} Langue maternelle}
}

\end{document}
